Ancestry-informative regions were strongly differentiated among the 20 hybrid populations, but distant regions usually differentiated different populations. Similarity in ancestry patterns declined rapidly with genetic distance and was close to its background level beyond about 5~cM and between chromosomes. Strong genome-wide differentiation therefore does not reflect a single ancestry pattern repeated across loci (Figure~\ref{fig:decay}).

Sorted loci also showed no greater ancestry-profile similarity or within-population LD with neighbouring markers than matched unsorted loci. This result held for both LD-reduced regions and all unpruned SNPs, including SNPs within each focal locus's own LD cluster. Sorted loci were therefore not embedded in longer ancestry blocks than equally differentiated unsorted loci in comparable recombination contexts. Twenty-five sorted loci in unusually large or otherwise poorly matched LD blocks could not be assessed and remain possible exceptions (Figure~\ref{fig:neighbourhood}; Table~\ref{tab:unmatched}).

Together, the results are consistent with ancestry sorting occurring largely region by region. Comparisons with corrected neutral simulations are still needed before attributing this pattern specifically to selection.
